\begin{textsample}{POS Dim 4 – rr – Score 12.00 – rr/rr\_em\_16.txt}  \label{ex:f4_pos_065}
2.3. \textbf{migrantes} O Estado de Roraima é historicamente desenvolvido pelos constantes e sucessivos fluxos \textbf{migratórios} e imigratórios, pois ao longo de sua história, têm acolhido brasileiros de outras regiões, principalmente nordestinos, como também \textbf{migrantes} de diversas nacionalidades, tais como, \textbf{venezuelanos}, guianenses, haitianos, cubanos, dentre outras.
Desse modo, fica evidente as especificidades e complexidades que o \textbf{território} de Roraima construiu ao longo do tempo.
Vale ( 2006, p.256 ) caracteriza esses processos ao refletir sobre o impacto da \textbf{migração} por transferência de \textbf{populações} em Roraima. > Uma das evidências dos resultados negativos da transferência de \textbf{populações} refere -se ao elevado ritmo do aumento demográfico, desproporcional ao que os governos implantam em infra-estrutura, acarretando problemas já conhecidos pelas comunidades dos grandes centros urbanos.
Quando falamos especificamente do fluxo de \textbf{migrantes} nordestinos para região \textbf{amazônica} e em especial para Roraima, percebe -se que está condicionado aos estímulos econômicos e geopolíticos.
Souza e Nogueira ( 2013 ) os descrevem muito bem quando mencionam o auge da exploração da borracha, a abertura de estradas nas décadas de 1950 e 1960, a política de ocupação do Regime Militar e as ações balizadas pelo conceito de desenvolvimento sustentável.
Neste contexto, Vale ( 2006 ), confere ao processo de \textbf{migração} nordestina em Roraima um papel relevante na produção de novas territorialidades e novas formas de concepção do uso e domínio do \textbf{território}, que se caracterizam pelos valores tradicionais desse povo, associado a uma carga de influência das \textbf{populações} nativas.
Atualmente, o estado de Roraima passa por um outro forte processe \textbf{migratório} que é caracterizado da seguinte maneira por Lima e Fernandes ( 2019 ). > O estado de Roraima ( Brasil ) nos últimos anos tem recebido \textbf{migrantes} \textbf{venezuelanos}, haitianos e cubanos.
Os cubanos entram pela fronteira da República Cooperativista da Guiana.
Já os haitianos vêm de Manaus ( Estado Amazonas, Brasil ) e algumas famílias vem da Venezuela.
Em 2016, houve uma intensificação da vinda de \textbf{venezuelanos} \textbf{indígenas} e não-indígenas para Roraima, principalmente para a cidade de Boa Vista, \textbf{capital} do estado.
Nesse contexto, é de suma importância que as redes de ensino, considerem esta realidade, desenvolvam estratégias de acolhimento e promovam ações que visem à integração dessas \textbf{populações}.
Além disso, Lima e Fernandes ( 2019 ) ao citar Marcelo Naputano, sugerem a necessidade de se desenvolver : > [... ] programa de educação voltado para a mediação socioeducativa entre gestores, professores, estudantes e famílias de refugiados \textbf{venezuelanos} como alternativa para aprofundar a integração cultural, onde se possa vivenciar a educação como experiência de uma construção relacional cotidiana em que se crie vínculos em meio aos conflitos.
Portanto, considerando os princípios que embasam o Novo Ensino Médio e foram incorporados no DCRR, as escolas precisam estar atentas para incluir os aspectos da \textbf{migração} local em seus projetos pedagógicos, na perspectiva de acolhimento das \textbf{populações} \textbf{migrantes} presentes no \textbf{território} roraimense de forma que promova a integração cultural e uma educação integral.
Ainda, cabe a escola, compreender que ambiente educativo, há uma composição de sujeitos com diversas identidades, culturas e história de vida, que devem ser considerados para a promoção de uma educação inclusiva, que assegure os direitos de todos os estudantes, conforme estabelece as legislações específicas referentes ao atendimento para este público alvo.

% matched lemmas: amazônico, capital, indígena, migrante, migratório, migração, população, território, venezuelano
\end{textsample}
